\documentclass[11pt, a4paper]{article}
\usepackage[a4paper, margin=1.5cm, bottom=1.8cm]{geometry}
\usepackage[T1]{fontenc}
\usepackage[utf8]{inputenc}
\usepackage[italian]{babel}
\usepackage{amsmath, amssymb}
\usepackage{enumitem}
\usepackage{xcolor}
\usepackage{tcolorbox}
\usepackage{booktabs}
\usepackage{fancyhdr}
\usepackage{hyperref}

% Palette cromatica per Liceo Sportivo
\definecolor{primary}{RGB}{15, 76, 129}   % Blu sportivo
\definecolor{secondary}{RGB}{211, 84, 0}   % Arancio energico
\definecolor{accent}{RGB}{39, 174, 96}     % Verde
\definecolor{bglight}{RGB}{248, 249, 250}  % Grigio chiarissimo per i box

\pagestyle{fancy}
\fancyhf{}
\renewcommand{\headrulewidth}{0pt}
\cfoot{\small\color{gray} Liceo Scientifico ad Indirizzo Sportivo -- Laboratorio di Matematica: Teoria degli Insiemi}

\setlist[enumerate]{leftmargin=*, topsep=2pt, itemsep=2pt, parsep=0pt}
\setlist[itemize]{leftmargin=*, topsep=2pt, itemsep=2pt, parsep=0pt}

\tcbset{
  colback=bglight,
  colframe=primary,
  fonttitle=\bfseries\color{white},
  boxrule=0.8pt,
  arc=3pt,
  top=4pt, bottom=4pt, left=6pt, right=6pt
}

\begin{document}

% ==============================================================================
% SQUADRA ALPHA - FACCIATA 1
% ==============================================================================

\begin{tcolorbox}[colframe=primary, title=\Large\bfseries SFIDA A SQUADRE: DAI PROBLEMI AGLI INSIEMI -- SQUADRA ALPHA]
\small
\textbf{Classe:} I Liceo Scientifico Sportivo \hfill \textbf{Tipologia Scheda:} ALPHA (Squadre 1, 2, 3, 4) \\
\textbf{Nome Squadra:} \dotfill \hfill \textbf{Data:} \dotfill \\
\textbf{Componenti e Ruoli:} 1. \dotfill (Capitano) -- 2. \dotfill (Scriba) \\
3. \dotfill (Revisore) -- 4. \dotfill (Relatore) [-- 5. \dotfill (Analista)]
\end{tcolorbox}

\vspace{-2mm}
\section*{Parte 1: Risoluzione di Problemi Pratici}
\textit{Lavorate in gruppo per risolvere gli 8 problemi sottostanti. Utilizzate diagrammi, schemi, tabelle o ragionamenti intuitivi per giustificare le risposte.}

\subsection*{Categoria A: Problemi su Insiemi e Frequenze}
\begin{enumerate}
    \item \textbf{Frequentatori della Palestra e del Centro Atletico:} Un'indagine condotta su un gruppo di studenti del liceo sportivo rileva che:
    \begin{itemize}
        \item 50 studenti frequentano la palestra scolastica e, fra questi, 42 \textbf{non} frequentano il centro atletico comunale;
        \item 60 studenti frequentano il centro atletico comunale;
        \item 18 studenti non frequentano né la palestra scolastica né il centro atletico comunale.
    \end{itemize}
    \textbf{Quesito:} Quanti studenti costituiscono complessivamente il gruppo intervistato? Disegnate un diagramma a cerchi (Euler-Venn) per rappresentare visivamente la situazione e mostrare i calcoli.
    
    \vspace{1.5cm}

    \item \textbf{Calcio e Tennis a Scuola:} In una scuola di 400 studenti, 150 praticano il calcio, 100 praticano il tennis e 30 praticano entrambe le attività sportive.
    \textbf{Quesito:} Quanti sono gli studenti della scuola che non praticano né il calcio né il tennis?
    
    \vspace{1.2cm}
\end{enumerate}

\subsection*{Categoria B: Combinatoria e Conteggi}
\begin{enumerate}[resume]
    \item \textbf{Il Menu dell'Atleta:} Il menu della mensa dello sportivo comprende: 4 antipasti energetici, 3 primi piatti proteici e 3 secondi piatti. In quanti modi distinti un atleta può comporre il proprio pranzo scegliendo esattamente un antipasto, un primo e un secondo?
    
    \vspace{1cm}

    \item \textbf{Strette di Mano al Torneo:} Al termine di una competizione, 6 atleti si incontrano e ciascuno stringe la mano una e una sola volta a tutti gli altri. Quante strette di mano si scambieranno complessivamente in tutto? Spiegate la regola seguita.
    
    \vspace{1cm}
\end{enumerate}

\subsection*{Categoria C: Ragionamento e Deduzioni Logiche}
\begin{enumerate}[resume]
    \item \textbf{Allenamento e Palestra:} Stabilite se il seguente ragionamento è valido o errato, motivando la risposta:
    \begin{quote}
    \textit{«Se Paolo fa parte della squadra di pallavolo, allora si allena tutti i giorni. Se Paolo si allena tutti i giorni, allora si reca tutti i giorni in palestra. Paolo non si reca tutti i giorni in palestra. Quindi: Paolo non fa parte della squadra di pallavolo.»}
    \end{quote}
    
    \vspace{1cm}

    \item \textbf{Erbivori e Dinosauri:} Stabilite se il seguente ragionamento è valido o errato, spiegando l'eventuale errore:
    \begin{quote}
    \textit{«Alcuni dinosauri sono erbivori. Otto è un erbivoro. Quindi: Otto è un dinosauro.»}
    \end{quote}
    
    \vspace{1cm}
\end{enumerate}

\subsection*{Categoria D: Negazioni e Quantificatori Logici}
\begin{enumerate}[resume]
    \item \textbf{Gli Asini Volanti:} La negazione logica corretta dell'affermazione «\textit{Esiste un asino che vola}» è:
    \begin{quote}
    A) Esiste un asino che cammina \quad B) Nessun asino vola \quad C) Esistono più asini che volano \\
    D) Tutti gli asini camminano \quad E) Esiste un asino che non vola
    \end{quote}
    \textbf{Scelta e motivazione:} \dotfill

    \vspace{0.8cm}

    \item \textbf{Oggetti d'Avorio:} La negazione logica corretta dell'affermazione «\textit{Ogni oggetto d'avorio è bianco}» è:
    \begin{quote}
    A) Ogni oggetto d'avorio è nero \quad B) Nessun oggetto d'avorio è bianco \quad C) Esiste un oggetto d'avorio che è nero \\
    D) Ogni oggetto d'avorio non è bianco \quad E) Esiste un oggetto d me d'avorio che non è bianco
    \end{quote}
    \textbf{Scelta e motivazione:} \dotfill
\end{enumerate}

\clearpage

% ==============================================================================
% SQUADRA ALPHA - FACCIATA 2
% ==============================================================================

\begin{tcolorbox}[colframe=secondary, title=\Large\bfseries PARTE 2: LABORATORIO DI FORMALIZZAZIONE (SQUADRA ALPHA)]
\textbf{Obiettivo:} Senza consultare il libro di testo, usate gli esercizi risolti nella Facciata 1 per completare in autonomia le definizioni matematiche formali degli insiemi.
\end{tcolorbox}

\vspace{-1mm}

\section*{1. Concetto di INTERSEZIONE ($\cap$)}
\begin{itemize}
    \item \textbf{Completate la definizione:} Dati due insiemi $A$ e $B$, l'insieme \textbf{Intersezione} (simbolo $A \cap B$) è l'insieme formato da tutti gli elementi che appartengono ad $A$ \textbf{\dots\dots\dots\dots} ad $B$ (\textit{inserire la congiunzione logica corretta: "e" oppure "o"}).
    \item \textbf{Collegamento all'Esercizio 2:} Nell'Esercizio 2 sul Calcio ($C$) e Tennis ($T$), quale numero rappresenta il numero di elementi dell'intersezione $C \cap T$? Risposta: \dotfill
\end{itemize}

\vspace{-1mm}

\section*{2. Concetto di UNIONE ($\cup$)}
\begin{itemize}
    \item \textbf{Completate la definizione:} Dati due insiemi $A$ e $B$, l'insieme \textbf{Unione} (simbolo $A \cup B$) è l'insieme formato da tutti gli elementi che appartengono ad $A$ \textbf{\dots\dots\dots\dots} ad $B$ (o a entrambi).
    \item \textbf{Collegamento all'Esercizio 2:} Quanti studenti praticano almeno uno tra il calcio e il tennis (cioè appartengono all'unione $C \cup T$)? Mostrate come si calcola questo numero partendo dai dati dell'esercizio:
    \begin{quote}
    $|C \cup T| = |C| + |T| - |C \cap T| = $ \dots\dots\dots\dots\dots\dots\dots\dots\dots\dots\dots\dots
    \end{quote}
\end{itemize}

\vspace{-1mm}

\section*{3. Concetto di DIFFERENZA tra Insiemi ($A \setminus B$)}
\begin{itemize}
    \item \textbf{Completate la definizione:} Dati due insiemi $A$ e $B$, l'insieme \textbf{Differenza} $A \setminus B$ (o $A - B$) è l'insieme formato dagli elementi che appartengono ad $A$ ma \textbf{\dots\dots\dots\dots\dots\dots\dots\dots} appartengono a $B$.
    \item \textbf{Collegamento all'Esercizio 1:} Nell'Esercizio 1, quanti studenti appartengono all'insieme "Frequentano la palestra ma NOT il centro atletico"? Questo insieme corrisponde alla differenza tra quali insiemi? \dotfill
    \item \textbf{Riflessione Logica:} La differenza $A \setminus B$ è uguale alla differenza $B \setminus A$? Spiegate portando un esempio pratico.
    
    \vspace{0.8cm}
\end{itemize}

\vspace{-1mm}

\section*{4. Concetto di SOTTOINSIEME ($A \subseteq B$)}
\begin{itemize}
    \item \textbf{Completate la definizione:} Si dice che un insieme $A$ è \textbf{sottoinsieme} di $B$ ($A \subseteq B$) se \textbf{\dots\dots\dots\dots\dots\dots} elemento che appartiene ad $A$ appartiene anche a $B$.
    \item \textbf{Collegamento all'Esercizio 6:} Se $D$ è l'insieme dei dinosauri ed $E$ è l'insieme degli erbivori, l'affermazione «\textit{Alcuni dinosauri sono erbivori}» significa che $D$ è sottoinsieme di $E$? Perché l'ipotesi che "Otto è un erbivoro" non ci permette di concludere con certezza che Otto sia un dinosauro?
    
    \vspace{0.8cm}
\end{itemize}

\vspace{-1mm}

\section*{5. Concetto di INSIEME VUOTO ($\emptyset$)}
\begin{itemize}
    \item \textbf{Completate la definizione:} L'insieme \textbf{Vuoto} è un insieme speciale che contiene \textbf{\dots\dots\dots\dots\dots\dots} elementi. Si indica con il simbolo $\emptyset$ oppure $\{ \}$.
    \item Se due insiemi $A$ e $B$ non hanno alcun elemento in comune, la loro intersezione è: $A \cap B = $ \dotfill
\end{itemize}

\vspace{-1mm}

\section*{6. Sintesi Grafica della Squadra}
Rappresentate con i diagrammi di Euler-Venn le seguenti operazioni, tratteggiando o colorando la regione corrispondente:

\begin{center}
\small
\begin{tabular}{ccccc}
\textbf{Intersezione } $A \cap B$ & \textbf{Unione } $A \cup B$ & \textbf{Differenza } $A \setminus B$ & \textbf{Differenza } $B \setminus A$ & \textbf{Sottoinsieme } $A \subseteq B$ \\[4pt]
\begin{minipage}{3cm}
\centering
\setlength{\unitlength}{0.6mm}
\begin{picture}(40,30)
\put(15,15){\circle{22}}
\put(25,15){\circle{22}}
\put(8,15){$A$}
\put(28,15){$B$}
\end{picture}
\end{minipage} &
\begin{minipage}{3cm}
\centering
\setlength{\unitlength}{0.6mm}
\begin{picture}(40,30)
\put(15,15){\circle{22}}
\put(25,15){\circle{22}}
\put(8,15){$A$}
\put(28,15){$B$}
\end{picture}
\end{minipage} &
\begin{minipage}{3cm}
\centering
\setlength{\unitlength}{0.6mm}
\begin{picture}(40,30)
\put(15,15){\circle{22}}
\put(25,15){\circle{22}}
\put(8,15){$A$}
\put(28,15){$B$}
\end{picture}
\end{minipage} &
\begin{minipage}{3cm}
\centering
\setlength{\unitlength}{0.6mm}
\begin{picture}(40,30)
\put(15,15){\circle{22}}
\put(25,15){\circle{22}}
\put(8,15){$A$}
\put(28,15){$B$}
\end{picture}
\end{minipage} &
\begin{minipage}{3cm}
\centering
\setlength{\unitlength}{0.6mm}
\begin{picture}(40,30)
\put(20,15){\circle{26}}
\put(20,15){\circle{14}}
\put(20,13){\tiny $A$}
\put(20,24){\tiny $B$}
\end{picture}
\end{minipage}
\end{tabular}
\end{center}

\clearpage

% ==============================================================================
% SQUADRA BETA - FACCIATA 1
% ==============================================================================

\begin{tcolorbox}[colframe=primary, title=\Large\bfseries SFIDA A SQUADRE: DAI PROBLEMI AGLI INSIEMI -- SQUADRA BETA]
\small
\textbf{Classe:} I Liceo Scientifico Sportivo \hfill \textbf{Tipologia Scheda:} BETA (Squadre 5, 6, 7, 8) \\
\textbf{Nome Squadra:} \dotfill \hfill \textbf{Data:} \dotfill \\
\textbf{Componenti e Ruoli:} 1. \dotfill (Capitano) -- 2. \dotfill (Scriba) \\
3. \dotfill (Revisore) -- 4. \dotfill (Relatore) [-- 5. \dotfill (Analista)]
\end{tcolorbox}

\vspace{-2mm}
\section*{Parte 1: Risoluzione di Problemi Pratici}
\textit{Lavorate in gruppo per risolvere gli 8 problemi sottostanti. Utilizzate diagrammi, schemi, tabelle o ragionamenti intuitivi per giustificare le risposte.}

\subsection*{Categoria A: Problemi su Insiemi e Frequenze}
\begin{enumerate}
    \item \textbf{Indagine sulle Lingue negli Atleti:} Da un'indagine sulla conoscenza delle lingue straniere effettuata su un gruppo di 500 atleti italiani, è emerso che:
    \begin{itemize}
        \item il 60\% conosce l'inglese;
        \item il 30\% conosce il francese;
        \item il 20\% non conosce né l'inglese né il francese.
    \end{itemize}
    \textbf{Quesito:} Quanti atleti conoscono sia l'inglese sia il francese? Quanti conoscono soltanto l'inglese? E quanti soltanto il francese? Disegnate il diagramma di Venn corrispondente.
    
    \vspace{1.5cm}

    \item \textbf{Inchiesta Polisportiva a Tre Attività:} Un'indagine condotta in un liceo sportivo rileva la pratica di tre discipline: Calcio ($C$), Nuoto ($N$) e Tennis ($T$). I dati sono:
    \begin{quote}
    a. 200 praticano calcio; \quad b. 140 nuoto; \quad c. 110 tennis; \quad d. 30 nuoto e tennis; \\
    e. 150 praticano calcio ma NON tennis né nuoto; \quad f. 20 praticano calcio e nuoto ma NON tennis; \\
    g. 10 praticano calcio, nuoto e tennis; \quad h. 40 non praticano nessuno dei tre sport.
    \end{quote}
    \textbf{Quesito:} Quanti sono complessivamente gli studenti della scuola?
    
    \vspace{1.2cm}
\end{enumerate}

\subsection*{Categoria B: Combinatoria e Conteggi}
\begin{enumerate}[resume]
    \item \textbf{Password degli Armadietti:} La serratura elettronica di un armadietto dello spogliatoio richiede una password di 4 caratteri.
    \begin{enumerate}
        \item Quante password diverse si possono generare usando unicamente le 10 cifre numeriche ($0-9$)?
        \item Quante password diverse si possono generare combinando cifre e 26 lettere dell'alfabeto?
    \end{enumerate}
    
    \vspace{1cm}

    \item \textbf{Il Guardaroba dello Sportivo:} Michele sceglie la propria tenuta per l'allenamento attingendo da: 4 paia di scarpe da ginnastica, 3 pantaloni di tuta, 5 magliette tecniche. Inoltre può scegliere se indossare o meno uno dei 3 cappellini da corsa che possiede. In quanti modi possibili ed elencabili può vestirsi?
    
    \vspace{1cm}
\end{enumerate}

\subsection*{Categoria C: Ragionamento e Deduzioni Logiche}
\begin{enumerate}[resume]
    \item \textbf{La Festa di Compleanno:} Stabilite se la seguente deduzione è valida o non valida, spiegando il motivo:
    \begin{quote}
    \textit{«Se Maria viene alla mia festa di compleanno, allora alla festa non viene Paolo. Se alla festa viene Marco, allora viene anche Paolo. Maria viene alla mia festa di compleanno. Quindi: Alla mia festa di compleanno non viene Marco.»}
    \end{quote}
    
    \vspace{1cm}

    \item \textbf{Concorso Viaggio in Messico:} Stabilite se la seguente deduzione è valida o non valida, spiegando il motivo:
    \begin{quote}
    \textit{«Se appartieni alle prime 100 persone che prenotano il viaggio, allora partecipi al concorso. Se partecipi al concorso, allora puoi vincere il viaggio gratis. Hai vinto il viaggio gratis. Quindi: Appartieni alle prime 100 persone che hanno prenotato il viaggio.»}
    \end{quote}
    
    \vspace{1cm}
\end{enumerate}

\subsection*{Categoria D: Negazioni e Quantificatori Logici}
\begin{enumerate}[resume]
    \item \textbf{Superamento Prova di Atletica:} Qual è la negazione logica corretta della frase «\textit{Tutti gli studenti della classe hanno superato la prova di atletica}»?
    \begin{quote}
    A) Nessuno studente ha superato la prova \quad B) Tutti gli studenti hanno fallito la prova \\
    C) Esiste almeno uno studente che non ha superato la prova \quad D) La maggior parte non l'ha superata
    \end{quote}
    \textbf{Scelta e motivazione:} \dotfill

    \vspace{0.8cm}

    \item \textbf{Calciatori a Piedi Scalzi:} Qual è la negazione logica corretta della frase «\textit{Esiste almeno un calciatore che gioca a piedi scalzi}»?
    \begin{quote}
    A) Nessun calciatore gioca a piedi scalzi \quad B) Esiste un calciatore che gioca con le scarpe \\
    C) Tutti i calciatori giocano a piedi scalzi \quad D) Alcuni calciatori non giocano a piedi scalzi
    \end{quote}
    \textbf{Scelta e motivazione:} \dotfill
\end{enumerate}

\clearpage

% ==============================================================================
% SQUADRA BETA - FACCIATA 2
% ==============================================================================

\begin{tcolorbox}[colframe=secondary, title=\Large\bfseries PARTE 2: LABORATORIO DI FORMALIZZAZIONE (SQUADRA BETA)]
\textbf{Obiettivo:} Senza consultare il libro di testo, usate gli esercizi risolti nella Facciata 1 per completare in autonomia le definizioni matematiche formali degli insiemi.
\end{tcolorbox}

\vspace{-1mm}

\section*{1. Concetto di INTERSEZIONE ($\cap$)}
\begin{itemize}
    \item \textbf{Completate la definizione:} Dati due insiemi $A$ e $B$, l'insieme \textbf{Intersezione} (simbolo $A \cap B$) è l'insieme formato da tutti gli elementi che appartengono ad $A$ \textbf{\dots\dots\dots\dots} ad $B$ (\textit{inserire la congiunzione logica corretta: "e" oppure "o"}).
    \item \textbf{Collegamento all'Esercizio 1:} Nell'Esercizio 1 (Inglese $I$ e Francese $F$), quanti atleti appartengono all'intersezione $I \cap F$? Come corrisponde questa percentuale all'area sovrapposta dei cerchi? \dotfill
\end{itemize}

\vspace{-1mm}

\section*{2. Concetto di UNIONE ($\cup$)}
\begin{itemize}
    \item \textbf{Completate la definizione:} Dati due insiemi $A$ e $B$, l'insieme \textbf{Unione} (simbolo $A \cup B$) è l'insieme formato da tutti gli elementi che appartengono ad $A$ \textbf{\dots\dots\dots\dots} ad $B$ (o a entrambi).
    \item \textbf{Collegamento all'Esercizio 1:} Qual è la percentuale di atleti che conoscono \textbf{almeno una} tra le due lingue (cioè $I \cup F$)? Se sommiamo $60\% + 30\% = 90\%$, perché otteniamo un valore maggiore? Cosa dobbiamo sottrarre?
    \begin{quote}
    $\% (I \cup F) = \% I + \% F - \% (I \cap F) = $ \dots\dots\dots\dots\dots\dots\dots\dots\dots\dots\dots\dots
    \end{quote}
\end{itemize}

\vspace{-1mm}

\section*{3. Concetto di DIFFERENZA tra Insiemi ($A \setminus B$)}
\begin{itemize}
    \item \textbf{Completate la definizione:} Dati due insiemi $A$ e $B$, l'insieme \textbf{Differenza} $A \setminus B$ (o $A - B$) è l'insieme formato dagli elementi che appartengono ad $A$ ma \textbf{\dots\dots\dots\dots\dots\dots\dots\dots} appartengono a $B$.
    \item \textbf{Collegamento all'Esercizio 2:} Nell'Esercizio 2, quanti studenti praticano il calcio ma non praticano né il nuoto né il tennis? Esprimete questo insieme usando le differenze.
    \item \textbf{Riflessione Logica:} La differenza tra l'insieme di chi parla Inglese e chi parla Francese ($I \setminus F$) dà lo stesso numero rispetto a $F \setminus I$? Rispondete con i numeri dell'Esercizio 1.
    
    \vspace{0.8cm}
\end{itemize}

\vspace{-1mm}

\section*{4. Concetto di SOTTOINSIEME ($A \subseteq B$)}
\begin{itemize}
    \item \textbf{Completate la definizione:} Si dice che un insieme $A$ è \textbf{sottoinsieme} di $B$ ($A \subseteq B$) se \textbf{\dots\dots\dots\dots\dots\dots} elemento che appartiene ad $A$ appartiene anche a $B$.
    \item \textbf{Collegamento all'Esercizio 6:} Nell'Esercizio 6 (prenotazione viaggio e concorso), essere tra le prime 100 persone ($P$) significa partecipare al concorso ($C$). $P$ è sottoinsieme di $C$ ($P \subseteq C$)? Perché vincere il concorso non garantisce di essere tra le prime 100 persone prenotate?
    
    \vspace{0.8cm}
\end{itemize}

\vspace{-1mm}

\section*{5. Concetto di INSIEME VUOTO ($\emptyset$)}
\begin{itemize}
    \item \textbf{Completate la definizione:} L'insieme \textbf{Vuoto} è un insieme speciale che contiene \textbf{\dots\dots\dots\dots\dots\dots} elementi. Si indica con il simbolo $\emptyset$ oppure $\{ \}$.
    \item Se due sport non hanno alcun praticante in comune, la loro intersezione è l'insieme \dotfill
\end{itemize}

\vspace{-1mm}

\section*{6. Sintesi Grafica della Squadra}
Rappresentate con i diagrammi di Euler-Venn le seguenti operazioni, tratteggiando o colorando la regione corrispondente:

\begin{center}
\small
\begin{tabular}{ccccc}
\textbf{Intersezione } $A \cap B$ & \textbf{Unione } $A \cup B$ & \textbf{Differenza } $A \setminus B$ & \textbf{Differenza } $B \setminus A$ & \textbf{Sottoinsieme } $A \subseteq B$ \\[4pt]
\begin{minipage}{3cm}
\centering
\setlength{\unitlength}{0.6mm}
\begin{picture}(40,30)
\put(15,15){\circle{22}}
\put(25,15){\circle{22}}
\put(8,15){$A$}
\put(28,15){$B$}
\end{picture}
\end{minipage} &
\begin{minipage}{3cm}
\centering
\setlength{\unitlength}{0.6mm}
\begin{picture}(40,30)
\put(15,15){\circle{22}}
\put(25,15){\circle{22}}
\put(8,15){$A$}
\put(28,15){$B$}
\end{picture}
\end{minipage} &
\begin{minipage}{3cm}
\centering
\setlength{\unitlength}{0.6mm}
\begin{picture}(40,30)
\put(15,15){\circle{22}}
\put(25,15){\circle{22}}
\put(8,15){$A$}
\put(28,15){$B$}
\end{picture}
\end{minipage} &
\begin{minipage}{3cm}
\centering
\setlength{\unitlength}{0.6mm}
\begin{picture}(40,30)
\put(15,15){\circle{22}}
\put(25,15){\circle{22}}
\put(8,15){$A$}
\put(28,15){$B$}
\end{picture}
\end{minipage} &
\begin{minipage}{3cm}
\centering
\setlength{\unitlength}{0.6mm}
\begin{picture}(40,30)
\put(20,15){\circle{26}}
\put(20,15){\circle{14}}
\put(20,13){\tiny $A$}
\put(20,24){\tiny $B$}
\end{picture}
\end{minipage}
\end{tabular}
\end{center}

\end{document}